% FILE: 02_lineage_and_context.tex
% Project: comparisons
% Role: cross-system-capability-comparison-corpus
% Protocol: v30.1.1

\section{Lineage and Context}
\label{sec:comparisons-lineage}

\subsection{2016 Language-Model Lesson}
\label{sec:comparisons-lineage-lm}

The thesis lineage begins with a 2016 language-model experiment whose finding was:
\emph{local text imitation does not conserve consequence}. A system that predicts
the next token of a process description does not thereby enforce that the described
process is lawful. The prediction surface and the process-execution surface are
categorically different. This lesson, while originating in language modeling,
generalizes directly to the process-mining comparison problem: a library that
correctly computes a fitness score does not thereby guarantee that the fitness
score is within a lawful range, or that the event log feeding the computation was
structurally valid, or that the error path is machine-readable.

The comparisons corpus is not a language-model artifact. It is the direct
consequence of asking, twelve years after that experiment: \emph{what would it
mean to enforce process-execution lawfulness at the compiler level rather than at
the runtime level?} The three-stack comparison is the answer applied to the
process-mining domain.

\subsection{Enterprise Process Gap}
\label{sec:comparisons-lineage-enterprise}

Between the 2016 experiment and the current research program, the enterprise
process gap widened. Process mining tools grew in capability but not in
enforcement rigor. PM4Py added OCEL 2.0 support, POWL discovery, and multi-system
streaming capabilities, while retaining the same fundamental architectural
decision: validation is the caller's responsibility, not the library's. A process
analyst who feeds an empty event log to PM4Py's inductive miner receives an empty
process tree without an error, not a rejection. A conformance analyst who receives
a fitness score of $0.97$ cannot determine from the score alone whether the
underlying log was OCEL-admitted or XES-admitted, or whether the score is bounded.

The gap the comparisons corpus documents is therefore not a gap in algorithmic
coverage alone --- wasm4pm is indeed missing Inductive Miner, alignment
conformance, and POWL discovery (\texttt{ALGORITHM\_COVERAGE\_MATRIX.md}). The
deeper gap is in \emph{enforcement architecture}: the absence of a layer that
makes invalid inputs and out-of-range outputs structurally impossible.

\subsection{Relationship to the Chatman Equation}
\label{sec:comparisons-lineage-chatman-eq}

The Chatman Equation, $A = \mu(O^*)$, frames the dissertation: an \emph{act}
$A$ is the image under the manufacturing map $\mu$ of an \emph{object} $O^*$.
In the process-mining context, a conformance verdict is an act. The comparisons
corpus establishes that the manufacturing map $\mu$ is not lawful unless:

\begin{enumerate}
  \item The input object $O^*$ (an event log) has been admitted through a named-law
    admission gate (\texttt{Evidence<T, Admitted, W>}).
  \item The output act $A$ (a conformance metric) is bounded by the type system
    (\texttt{Metric<KIND, NUM, DEN>} with \texttt{Require<\{NUM <= DEN\}>:
    IsTrue}).
  \item The refusal path, when admission fails, encodes which structural law was
    violated (\texttt{Refusal<R, W>} with a specific named-law type \texttt{R}).
\end{enumerate}

The thesis lineage from the Chatman Equation to the comparisons corpus is thus:
the equation demands that every act be traceable to a lawfully manufactured
object. The comparisons corpus operationalizes that demand across the three-stack
comparison. The connection from the Chatman Equation to these comparison tables is
a matter of architectural interpretation (AUTHOR INTERPRETATION) --- no document
in the comparisons directory itself names the Chatman Equation by that label.

\subsection{Relationship to Receipts and Replay}
\label{sec:comparisons-lineage-receipts}

The receipts and replay layer of the research program is the cryptographic
evidence chain that seals each manufacturing act. The comparisons corpus is
directly upstream of two canonical receipts in the receipt registry:

\begin{itemize}
  \item \textbf{PM4PY\_ORACLE\_RECEIPT} --- maps every PM4Py API function to its
    wasm4pm equivalent or GAP status, classifying gap severity as CRITICAL, HIGH,
    MEDIUM, or LOW. The comparison corpus provides the evidentiary basis for
    those severity classifications.
  \item \textbf{WASM4PM\_GAP\_RECEIPT} --- certifies that eight identified gaps
    have severity, priority, and wasm4pm-compat bridging paths documented. The
    algorithm coverage matrix in the comparisons corpus is the source of those
    gap entries.
\end{itemize}

The lattice compliance model in \texttt{lattice\_vs\_trace\_compliance\_checking.md}
also directly prefigures the receipt structure: the BLAKE3 hash chain over witness
states $H_k = \text{BLAKE3}(H_{k-1} \| e_k \| W_k)$ is the mathematical basis
for the cryptographic attestation that conformance receipts carry. The
comparisons corpus thus supplies both the gap inventory and the formal model for
the receipt layer.
